\documentclass{subfiles}
\begin{document}
    The first problem we introduce is that of finding frequents element in a stream.
        More formally, given a stream \(\sigma = \iprod{\sigma_{1}, \ldots, \sigma_{m}}\) and \(0 < \varepsilon < 1\), 
        we are asked to find the set \(I_{\varepsilon}(\sigma)\) that contains all 
        the elements \(j\) in \(\sigma\) such that \(F_{\sigma}[j] \ge \varepsilon m\).
        One of such elements is called \(\varepsilon\)-frequent, or \(\varepsilon\)-heavy hitter in \(\sigma\).

    When \(\varepsilon = \sfrac{1}{2}\), an \(\varepsilon\)-frequent token is also called a \emph{majority} item.
        A variant of the above problem, is therefore that of deciding whether a stream has 
        a majority item.

    Although the two problems appear to be easy, even in the vanilla model,
        the space required to solve them cannot be exactly sub-linear, 
        as proved by the following theorem.

    \begin{theorem*}
        Any deterministic streaming algorithm that solves the majority problem exactly on
        streams with \(m\) tokens over the universe \(U\), where \(m \le \sfrac{n}{2}\),
        must use \(\ofW{m log(n/m)}\) bits of storage, even in the vanilla streaming model.
    \end{theorem*}
    \begin{proof}
        Our proof is more of a sketch, some more rigorous proof are easely found online. 
        Let \(s\) be the number of bits used by an algorithm to solved the majority problem,
        then at any given time the algorithm can be in \(2^{s}\) different states.
        On the other hand, the number of possible different frequency vector is 
        \[
            \binom{n + \sfrac{m}{2} - 1}{n -1} =
            \binom{n + \sfrac{m}{2} - 1}{\sfrac{m}{2}} 
            \ge 2^{(\sfrac{m}{2})\log(2\sfrac{(n + \sfrac{m}{2} - 1)}{m}) }
        \]
        Hence, when \(s < (\sfrac{m}{2})\log(2\sfrac{(n + \sfrac{m}{2} - 1)}{m}) \) 
        there must exist two sequences \(t_{1} \ne t_{1}'\) such that the algorithm would 
        be in the same state when processing either of the two sequences.
        Then, ther exist a sequence \(t_{2}\) such that \(t_{1} \circ t_{2}\)
        contains a majority item \(j\), while \(j\) is not a majority item in \(t_{1}' \circ t_{2}\).
        It follows that the algorithm would report the same answer for both cases, 
        since it is the same state after processing either of sequences.
        But it is incorrect, since \(j\) is a majority item for \(t_{1} \circ t_{2}\)
        and not for \(t_{1}' \circ t_{2}\).
    \end{proof}

    Despite the result given by the above theorem, 
        it can be proved that the majority proble can be solved using \(\ofO{\log (m + n}\) bits.
        The idea underlying the algorithm we present is the following: suppose that,
        given a sequence of numbers, we repeatedly delete any pair of adjacent and distinct numbers,
        until either all remaining numbers are equal or the sequence has become empty. 
        Then the remaining number, if any, is the only candidate for a majority item.
        The algorithm in \Cref{alg:majority} works as follow: given a threshold \(\varepsilon\),
        it maintains a set \(I\) that is guaranteed to have all \(\varepsilon\)-frequent 
        elements, where \(\abs{I} < \sfrac{1}{\varepsilon}\). For each element in
        the set, it maintains a counter \(c(j)\) to keep track of the occurrences of \(j\).

    \subfile{../TikZ/Algorithm 1 - Majority}

\end{document}
